\section{特别程序}
\begin{enumerate}
  \item \zd{牢记特别程序是非诉程序，不解决权利义务纠纷}。\al{调解、辩论、质证、开庭、证据 \cuo}。
  \item \zd{选民资格}、宣告失踪/死亡、无民/限民认定、认定财产无主、确认调解协议效力、实现担保物权。
  \item 共同特点。
        \begin{enumerate}
          \item 原则上\zd{独任制}；选民资格和重大疑难案件合议庭；\zd{没有人民陪审员}。
          \item 一律\zd{一审终审}；判决\zd{作出生效}；且不适用三撤、再审制度（提异议 \hz 诉讼）。
          \item \zd{有错提异议}，适用特别程序作出新判决，裁定撤销原判决裁定。
          \item 立案之日 \hz 公告之日起，\zd{30 天内审结}；选民资格 \yjt 选举开始前。
        \end{enumerate}
  \item 选民资格。申诉前置。
  \item 宣告失踪/死亡。
        \begin{enumerate}
          \item 变更财产代管人。\yi 自己 \yjt 参照适用特别程序。\er 其他权利人 \yjt 以代管人为被告起诉，走普通程序。
        \end{enumerate}
  \item \zd{确认（人民调解委员会）调解协议效力/司法确认}。\al{确认和解协议效力 \cuo}。
        \begin{enumerate}
          \item 生效之日起 \zd{30 天内共同}提出申请，\zd{基层和中院}（本身属于中院管辖）都可以。
                \begin{enumerate}
                  \item 自行调解 \yjt \zd{当事人}住所地、\zd{标的物}所在地、\zd{调解组织}所在地法院。
                  \item 法院组织的\zd{诉前}调解 \yjt 谁邀请找谁。
                \end{enumerate}
          \item 不予受理。
                \begin{enumerate}
                  \item 不属于法院受理范围，不属于收到申请的法院管辖。
                  \item 涉及婚姻、亲子、收养\zd{身份关系}的。
                  \item 涉及\zd{物权、知识产权}确权。水上加水。\al{甲乙达成调解协议，房屋归甲}。
                  \item 涉及\zd{适用其他特别程序}。
                \end{enumerate}
          \item 肯定。确认调解协议效力 \yjt 给救济 \yjt \zd{提异议}。\yi 当事人 15 天。\er 利害人知道 \hz 应当知道 6 个月。
          \item 否定。驳回 \yjt 走和解、调解、仲裁、（\zd{就调解协议}）诉讼。
        \end{enumerate}
  \item \zd{实现担保物权}。为了配合物权法的实施，不走诉讼程序就能作出一个强制执行的法律文书出来。
        \begin{enumerate}
          \item 申请人。担保物权人（抵押、质押、留置权人）、其他权利人（质押 \hz 留置中的债务人）。
          \item 被申请人。执行标的所有人。
          \item 由担保\zd{财产所在地 \hz 担保物权登记地的基层法院}管辖。海事的归海事法院。
          \item 独任制。\zd{标的额超了的（视为重大疑难案件）应当组成合议庭审查}。
          \item 抵押贷款合同存在\zd{仲裁协议，不影响}当事人通过\zd{实现担保物权等特别程序}去实现权利。
                \begin{enumerate}
                  \item 申请 \yjt 就\zd{没有争议的部分}，裁定准许\zd{拍卖、变更 \yjt 执行}。
                  \item 申请 \yjt 有争议，驳回 \yjt 有纠纷 \yjt 走和解、调解、仲裁、\scx{诉讼}。
                \end{enumerate}
          \item 审查、救济同确认调解协议效力。
        \end{enumerate}
\end{enumerate}
